\section{问题四模型的建立与求解}

\subsection{建模思路与电价假设}
问题四将电价从“每天相同”放宽为逐日实时波动（附件 4，365 天$\times$144 点），
要求在波动电价下重新建立问题二与问题三条件下的购电策略模型。
本文采用电力市场的通行惯例并显式声明：

\textbf{假设（日前电价可知）}：每日 0:00 制定计划时，当日 144 点电价曲线已由日前市场公布
（可知但逐日不同）；紧急购电按交易时刻实际电价的 5 倍结算，
调整购电按附件 4 实际价格结算（超额 1.5 倍、违约 0.5 倍，同问题三）。
日前电价未知的情形作为推广方向讨论。

在此假设下，问题二、三已建立的策略框架（预测$\to$日前 LP 计划$\to$滚动调整$\to$日内实时调度）
\textbf{结构完全复用}，唯一变化是所有价格项（计划目标、紧急费率、调整结算）
由固定曲线 $c_t$ 换成当日曲线 $c_{d,t}$。

\subsection{附件 4 电价特征分析}
对附件 4 全年电价作统计，结果如表~\ref{tab:q4-price} 所示。

\begin{table}[!htbp]
\caption{附件 4（波动电价）与附件 1（固定电价）的统计特征对比}\label{tab:q4-price}
\centering
\small
\begin{tabular}{lcc}
\toprule
指标 & 附件 4（波动） & 附件 1（固定） \\
\midrule
全年均值/（元$\cdot$kWh$^{-1}$） & 0.7662 & 0.7662 \\
取值范围/（元$\cdot$kWh$^{-1}$） & [0.0076,\ 1.7936] & [0.3713,\ 1.3952] \\
日均价范围/（元$\cdot$kWh$^{-1}$） & [0.5378,\ 0.9290] & 0.7662（每日相同） \\
日内价差（max$-$min）均值/（元$\cdot$kWh$^{-1}$） & 1.1073 & 1.0239 \\
\bottomrule
\end{tabular}
\end{table}

全年均值与附件 1 \textbf{完全相同}——附件 4 是“均值不变、只打散波动”的设计，
排除了“电价整体上移”的解释，费用变化全部来自结构性因素。
进一步考察电价与负载、光伏的日尺度相关结构（图~\ref{fig:q4-corr}）：

\textbf{（1）电价与负载强正相关}：corr（日均价，日负载电量）$=0.898$。
高负载日几乎必然是高电价日——能量必须在最贵的日子买入，无法跨日回避。

\textbf{（2）电价与光伏负相关}：corr（日均价，日光伏电量）$=-0.343$。
光伏大发时电价偏低、缺发时电价偏高，光伏替代购电的“保值”作用被进一步削弱。

\begin{figure}[!htbp]
\centering
\includegraphics[width=0.9\textwidth]{Q4_price_pv.pdf}
\caption{附件 4 日均价与日负载电量、日光伏电量的散点及回归线：
(a) 价—荷相关系数 $+0.898$；(b) 价—光相关系数 $-0.343$}
\label{fig:q4-corr}
\end{figure}

\subsection{波动电价下的购电策略模型}
\textbf{（1）Q4-2（对应问题二）。} w1 预测 $+$ 余量（负载 $\times$1.05、光伏 $\times$0.98）
$+$ 日前 LP（$S_{144}=S_0$ 日循环）$+$ 日内储能纠偏 $+$ 5 倍紧急购电；
日前 LP 目标为 $\min\sum_t c_{d,t}\,p_{d,t}\,\Delta t$，$c_{d,t}$ 为附件 4 第 $d$ 日曲线。
余量参数在附件 4 上复扫后仍最优（详见第九章）。

\textbf{（2）Q4-3（对应问题三）。} 标定光伏估计 $+$ 当日残差更新 $+$ 0:00 计划 LP
$+$ 6:00/12:00/18:00 双向调整 LP（$c_{d,t}$ 计入目标与安全阀）$+$ 日内纠偏；
结算同问题三：$\sum c_{d,t}p^{adj}_{d,t}\Delta t+\sum 0.5c_{d,t}|p^{adj}_{d,t}-p^{plan}_{d,t}|\Delta t+\sum 5c_{d,t}e_{d,t}\Delta t$。
余量参数在波动电价下重新寻优：$\gamma=1.02$ 保持最优，光伏余量由 $\delta=0.98$ 调整为 $\delta=1.00$
（波动电价下光伏边际价值下降，过度下压光伏预测反而抬高购电成本，见第九章表~\ref{tab:sens-q4}）；
调整 LP 沿用日末最低备用约束 $S_{144}\ge4000$\,kWh。

\subsection{求解结果}
两种策略在波动电价下的全年滚动仿真结果如表~\ref{tab:q4-total} 与表~\ref{tab:q4-comp} 所示。

\begin{table}[!htbp]
\caption{固定电价与波动电价下的全年总费用对比（报告期 334 天）}\label{tab:q4-total}
\centering
\small
\setlength{\tabcolsep}{5pt}
\begin{tabular}{lccc}
\toprule
策略 & 固定电价（附件 1） & 波动电价（附件 4） & 变化 \\
\midrule
Q4-2（问题二策略） & 14546503.97 元 & \textbf{15263635.75 元} & $+$717131.78（$+$4.93\%） \\
Q4-3（问题三策略 V4FR） & 13687490.69 元 & \textbf{14315636.38 元} & $+$628145.69（$+$4.59\%） \\
Q4-3 相对 Q4-2 & $-$5.91\% & \textbf{$-$6.21\%} & 优势扩大 \\
\bottomrule
\end{tabular}
\end{table}

\begin{table}[!htbp]
\caption{波动电价下的费用构成（单位：元）}\label{tab:q4-comp}
\centering
\small
\setlength{\tabcolsep}{5pt}
\begin{tabular}{lcccc}
\toprule
策略 & 计划/实付购电费 & 调整违约/溢价 & 紧急购电费 & 紧急购电量/天数 \\
\midrule
Q4-2 & 14336659.77 & --- & 926975.98 & 146798.3\,kWh / 105 天 \\
Q4-3 & 13568042.27 & 364362.21 & 383231.90 & 70819.5\,kWh / 199 天 \\
\bottomrule
\end{tabular}
\end{table}

四个指定日期在波动电价下的计算结果如表~\ref{tab:q4-days} 所示（各日期指定时段购电量与储能
充放电量明细见结果文件 result4-2.xlsx 与 result4-3.xlsx）。

\begin{table}[!htbp]
\caption{四个指定日期在波动电价下的购电结果}\label{tab:q4-days}
\centering
\small
\setlength{\tabcolsep}{4pt}
\begin{tabular}{lccccc}
\toprule
 & \multicolumn{2}{c}{Q4-2（问题二策略）} & \multicolumn{3}{c}{Q4-3（问题三策略）} \\
\cmidrule(lr){2-3}\cmidrule(lr){4-6}
日期 & 计划购电量/kWh & 购电费/元 & 计划$\to$调整/kWh & 计划费$\to$结算费/元 & 紧急购电 \\
\midrule
2025-03-20 & 79420.62 & 49650.82 & 69723.56$\to$70905.44 & 43137.99$\to$44365.27 & 无 \\
2025-06-21 & 38848.12 & 19143.71 & 35919.35$\to$36269.04 & 17318.23$\to$18024.19 & 无 \\
2025-09-23 & 71264.87 & 45711.78 & 66530.09$\to$69322.40 & 42714.01$\to$46352.32 & 合计 156.48\,kWh \\
2025-12-21 & 100592.96 & 76223.89 & 95937.48$\to$98112.10 & 71680.77$\to$74388.56 & 8:00--8:10，34.76\,kWh \\
\bottomrule
\end{tabular}
\end{table}

Q4-2 仅 12 月 21 日发生紧急购电（8:00--8:10，5.62\,kWh），其余三天无；
Q4-3 的 9 月 23 日紧急购电为三段合计（9:20--10:00 共 149.02\,kWh、18:30--18:40 共 1.08\,kWh、
20:30--20:40 共 6.38\,kWh）。

\subsection{结果分析：费用上升的来源与策略适应性}
\textbf{（1）费用上升的主因是“价—荷同峰”。} 全年均值不变的前提下，
高电价恰好落在高负载日：负载加权的购电均价由固定电价的 0.7955 元/kWh 升至 0.8190 元/kWh（$+$2.95\%），
构成费用上升的主要部分；在“每日 0:00 计划 $+$ 日循环”的框架下无法跨日转移负载以规避
（跨日套利的价值见下述敏感性实验）。

\textbf{（2）光伏的避险价值被削弱。} 价—光负相关（$-0.343$）意味着发电时段电价低、缺发时段电价高，
光伏替代购电所避免的单位成本下降；这也是 Q4-3 最优光伏余量由 0.98 变为 1.00 的原因。

\textbf{（3）问题三策略对价格波动的适应性更好。} V4FR 在波动电价下依然最优，
且优势由 5.91\% 扩大至 6.21\%（节省 94.8 万元）；
其上调购电量与当日日内价差正相关（0.194）——高价差日更多使用 1.5 倍上调通道规避 5 倍高价紧急购电。
尾部风险也显著压平：Q4-2 最贵单个 10\,min 紧急购电达 8514 元（7 月 3 日 20:40，949.34\,kWh，电价 1.7936 元）、
最贵一天 81112 元（6 月 1 日），紧急购电费的 92.5\% 落在电价 $>1.0$ 元的时段；
Q4-3 对应指标降为 5931 元、17792 元与 67.4\%。

\textbf{（4）波动电价激活了更多储能套利周转。} 全年储能充电量由固定电价下的
6653.2\,MWh（Q4-2）/6502.4\,MWh（Q4-3）升至 7123.4\,MWh（$+$7.1\%）/6840.1\,MWh（$+$5.2\%），
约每日 2.0 次等效循环；但日维度“日内价差—当日充电量”相关仅 $+0.18$，
说明周转提升主要来自价格形态的整体变化，容量与功率限额仍是主要约束，
与问题一中“扩容有正边际收益”的结论一致。

\textbf{（5）敏感性：跨日套利的价值。} 附件 4 日均价在 0.54$\sim$0.93 元之间波动，
跨日价差甚至大于日内价差。放开 Q4-2 的日循环约束（计划 LP 末端自由）重算，
总费用降至 15031398 元，再省 232238 元（$-$1.5\%）——跨日套利是波动电价下真实的价值来源；
但代价是储能频繁贴限（日末均值 2355\,kWh）、削弱次日备用，
实际运用中须在备用水平与套利收益之间权衡，可作为后续改进方向。

综上，波动电价通过“价—荷同峰”结构使两种策略费用分别上升 4.93\% 与 4.59\%，
问题三的“标定 $+$ 双向调整”框架在波动电价下保持了优势并压平了尾部风险，
模型结构无需改变即可量化价格波动的影响。
